\documentclass{article}
\usepackage{mdposter}
\posttitle{Growth of pattern-avoiding ascent sequences}
\postfield{Mathematics, enumerative combinatorics}

\begin{document}
\begin{mdposter}

\posterbox{name=c1, column=1, between=top and bottom}{%
\begin{mdcol}
\mdsec{The question}
An ascent sequence starts at 0, and each new entry is at most one more than the number of ascents so far.
How many ascent sequences of length $n$ avoid a family of patterns such as 0120, 0121 and 0122, and how fast
does that number $A(n)$ grow?
\mdgap
\mdsec{How we are approaching it}
\begin{itemize}
\item Join consecutive occurrences of the same value. These arcs form a repetition graph.
\item Bound its chromatic number, so that a sequence can be rebuilt from one occurrence of each value and a little extra data.
\item Compare with permutations for the upper bound, and append distinct larger entries to suitable prefixes for the lower bound.
\end{itemize}
\mdgap
\mdsec{What we have so far}
\begin{progress}
For any family of patterns $012\cdots(d{+}u)d$ that each use at least three values,
$\log A(n) = n\log n - 2n\log\log n + O(n)$.
\end{progress}
\begin{progress}
The generating functions are not D-finite, and the exponential one is entire of infinite order. With a pattern of fewer than three values, the generating function is rational.
\end{progress}
\end{mdcol}
}

\posterbox{name=c2, column=2, between=top and bottom}{%
\begin{mdcol}
\mdsec{Exhaustive counts}
\includegraphics[width=\linewidth]{ascent-counts.pdf}
\figcap{Figure 1.}{We counted by exhaustive search, up to $n = 13$, the sequences of each length that avoid 0120, 0121 and 0122 (blue) and all ascent sequences (grey). At $n = 13$, 44,616 of the 89,903,100 ascent sequences avoid the patterns.}
\mdgap
\mdsec{A repetition graph}
\includegraphics[width=\linewidth]{ascent-example.pdf}
\figcap{Figure 2.}{This avoiding sequence has length 30, and each arc joins consecutive occurrences of the same value.}
\end{mdcol}
}

\posterbox{name=c3, column=3, between=top and bottom}{%
\begin{mdcol}
\mdsec{In Mantis}
\mantis{ascent-sequences.jpg}{3,230}{The product-order theorem uses the unique repeated value in $R_{(d,u)}$.}{%
The same record corrects an earlier example: $\{(0,2), (2,2)\}$ had been claimed as an antichain, and the replacement is $\{(0,3), (2,1)\}$.\alsoonrecord Review caught an invalid weakening of an inequality.}
\mdgap
\mdsec{Open questions}
\begin{itemize}
\item The linear term of $\log A(n)$ is known for 0120, 0121 and 0122 avoided together. What is it for other families?
\item Which other families of sequences grow like $n\log n - 2n\log\log n$?
\end{itemize}
\end{mdcol}
}

\end{mdposter}
\end{document}
